% GENERATED FILE. Edit canonical lessons, then run npm run generate:books.
% canonical-source-hash: fnv1a64:37630a6c7d11fc4c
% canonical-lessons: TE-C232-jadi, TE-C232-basta, TE-C232-totti, TE-C232-pipa, TE-C232-maggu

\chapter{Jar, Sack, Tub, Barrel, Mug}
\label{ch:te-jar-sack-tub-barrel-mug}
\hlchaptermodality{232}

\begin{chapteropening}
\textbf{By the end of this chapter:} \emph{I can say a jar, a sack, a tub, a barrel and a mug.}

Complete the last lesson of chapter 232: I can say a jar, a sack, a tub, a barrel and a mug.
\end{chapteropening}

\section[\texorpdfstring{\te{జాడీ} (jāḍī)}{జాడీ (jāḍī)}]{\te{జాడీ} (jāḍī) — a jar (for pickle)}

\label{lesson:TE-C232-jadi}



\begin{quote}
Before the new one: say the Telugu for glue, then the Telugu for a ruler.
\end{quote}

\subsection*{You'll want to know: \te{జాడీ}}
\textbf{\te{జాడీ}} — \emph{jāḍī} — \textquotedblleft{}a jar (for pickle)\textquotedblright{}.

\begin{grammarlens}[title={What you've built}]
One more word for this chapter.
\end{grammarlens}

\subsection*{Your turn}
\begin{itemize}
\raggedright
  \item \emph{Say it:} \emph{jāḍī}
  \item \emph{Say it:} \emph{jāḍī}, once more
  \item \emph{Say it:} say \emph{jāḍī}
  \item \emph{Recall:} say the Telugu for glue, then the Telugu for a ruler, then say \emph{jāḍī} again
\end{itemize}

\begin{tcolorbox}[breakable,colback=teal!4,colframe=teal!35!black,title={Before you move on}]
What does \te{జాడీ} mean? (\textquotedblleft{}A jar (for pickle)\textquotedblright{}.) Say it once more.
\end{tcolorbox}
\section[\texorpdfstring{\te{బస్తా} (bastā)}{బస్తా (bastā)}]{\te{బస్తా} (bastā) — a sack (of grain)}

\label{lesson:TE-C232-basta}



\begin{quote}
Before the new one: say the Telugu for a ruler, then the Telugu for a jar.
\end{quote}

\subsection*{You'll want to know: \te{బస్తా}}
\textbf{\te{బస్తా}} — \emph{bastā} — \textquotedblleft{}a sack (of grain)\textquotedblright{}.

\begin{grammarlens}[title={What you've built}]
One more word for this chapter.
\end{grammarlens}

\subsection*{Your turn}
\begin{itemize}
\raggedright
  \item \emph{Say it:} \emph{bastā}
  \item \emph{Say it:} \emph{bastā}, once more
  \item \emph{Say it:} say \emph{bastā}
  \item \emph{Recall:} say the Telugu for a ruler, then the Telugu for a jar, then say \emph{bastā} again
\end{itemize}

\begin{tcolorbox}[breakable,colback=teal!4,colframe=teal!35!black,title={Before you move on}]
What does \te{బస్తా} mean? (\textquotedblleft{}A sack (of grain)\textquotedblright{}.) Say it once more.
\end{tcolorbox}
\section[\texorpdfstring{\te{తొట్టి} (toṭṭi)}{తొట్టి (toṭṭi)}]{\te{తొట్టి} (toṭṭi) — a tub, a trough}

\label{lesson:TE-C232-totti}



\begin{quote}
Before the new one: say the Telugu for a jar, then the Telugu for a sack.
\end{quote}

\subsection*{You'll want to know: \te{తొట్టి}}
\textbf{\te{తొట్టి}} — \emph{toṭṭi} — \textquotedblleft{}a tub, a trough\textquotedblright{}.

\begin{grammarlens}[title={What you've built}]
One more word for this chapter.
\end{grammarlens}

\subsection*{Your turn}
\begin{itemize}
\raggedright
  \item \emph{Say it:} \emph{toṭṭi}
  \item \emph{Say it:} \emph{toṭṭi}, once more
  \item \emph{Say it:} say \emph{toṭṭi}
  \item \emph{Recall:} say the Telugu for a jar, then the Telugu for a sack, then say \emph{toṭṭi} again
\end{itemize}

\begin{tcolorbox}[breakable,colback=teal!4,colframe=teal!35!black,title={Before you move on}]
What does \te{తొట్టి} mean? (\textquotedblleft{}A tub, a trough\textquotedblright{}.) Say it once more.
\end{tcolorbox}
\section[\texorpdfstring{\te{పీపా} (pīpā)}{పీపా (pīpā)}]{\te{పీపా} (pīpā) — a barrel}

\label{lesson:TE-C232-pipa}



\begin{quote}
Before the new one: say the Telugu for a sack, then the Telugu for a tub.
\end{quote}

\subsection*{You'll want to know: \te{పీపా}}
\textbf{\te{పీపా}} — \emph{pīpā} — \textquotedblleft{}a barrel\textquotedblright{}.

\begin{grammarlens}[title={What you've built}]
One more word for this chapter.
\end{grammarlens}

\subsection*{Your turn}
\begin{itemize}
\raggedright
  \item \emph{Say it:} \emph{pīpā}
  \item \emph{Say it:} \emph{pīpā}, once more
  \item \emph{Say it:} say \emph{pīpā}
  \item \emph{Recall:} say the Telugu for a sack, then the Telugu for a tub, then say \emph{pīpā} again
\end{itemize}

\begin{tcolorbox}[breakable,colback=teal!4,colframe=teal!35!black,title={Before you move on}]
What does \te{పీపా} mean? (\textquotedblleft{}A barrel\textquotedblright{}.) Say it once more.
\end{tcolorbox}
\section[\texorpdfstring{\te{మగ్గు} (maggu)}{మగ్గు (maggu)}]{\te{మగ్గు} (maggu) — a mug}

\label{lesson:TE-C232-maggu}



\begin{quote}
Before the new one: say the Telugu for a tub, then the Telugu for a barrel.
\end{quote}

\subsection*{You'll want to know: \te{మగ్గు}}
\textbf{\te{మగ్గు}} — \emph{maggu} — \textquotedblleft{}a mug\textquotedblright{}.

\begin{grammarlens}[title={What you've built}]
One more word for this chapter.
\end{grammarlens}

\subsection*{Your turn}
\begin{itemize}
\raggedright
  \item \emph{Say it:} \emph{maggu}
  \item \emph{Say it:} \emph{maggu}, once more
  \item \emph{Say it:} say \emph{maggu}
  \item \emph{Recall:} say the Telugu for a tub, then the Telugu for a barrel, then say \emph{maggu} again
\end{itemize}

\begin{tcolorbox}[breakable,colback=teal!4,colframe=teal!35!black,title={Before you move on}]
What does \te{మగ్గు} mean? (\textquotedblleft{}A mug\textquotedblright{}.) Say it once more.
\end{tcolorbox}
